\documentclass[17pt,compress]{beamer}
\usepackage{beamerthemesplit}
\mode<presentation>
{
  \usetheme{Warsaw}
  \setbeamercovered{transparent}
  \setbeamertemplate{navigation symbols}{}
}
% Taken from Fernando's slides.
\usepackage{ae,aecompl}
\usepackage[scaled=.95]{helvet}

\usepackage[english]{babel}
\usepackage[latin1]{inputenc}
\usepackage[T1]{fontenc}

\newcounter{saveenumi}
\newcommand{\seti}{\setcounter{saveenumi}{\value{enumi}}}
\newcommand{\conti}{\setcounter{enumi}{\value{saveenumi}}}


\definecolor{blue}{rgb}{0.16,0.32,0.75}
\setbeamercolor{structure}{fg=blue}
\author[FOSSEE]{}
\institute[IIT Bombay]{}
\date[]{}
% \setbeamercovered{transparent}

% theme split
\usepackage{verbatim}
\newenvironment{colorverbatim}[1][]%
{%
\color{blue}
\verbatim
}%
{%
\endverbatim
}%

\usepackage{mathpazo,courier,euler}
\usepackage{listings}
\lstset{language=sh,
    basicstyle=\ttfamily\bfseries,
  showstringspaces=false,
  keywordstyle=\color{black}\bfseries}

% logo
\logo{\includegraphics[height=1.30 cm]{3t-logo.pdf}}
\logo{\includegraphics[height=1.30 cm]{fossee-logo.png}

\hspace{7.5cm}
\includegraphics[scale=0.3]{fossee-logo.png}\\
\hspace{281pt}
\includegraphics[scale=0.80]{3t-logo.pdf}}

\begin{document}

\sffamily \bfseries
\title
[Matrices]
{Python - Matrices}
\author
[FOSSEE, IIT - Bombay]
{\small Talk to a Teacher\\{\color{blue}\url{http://spoken-tutorial.org}}\\National Mission on Education
 through ICT\\{\color{blue}\url{http://sakshat.ac.in}} \\[0.5cm]{\tiny Script by: Hardik ghaghada \\ Naration by: Hardik Ghaghada \\ 13 August 2015}}

% slide 1
\begin{frame}
   \titlepage
\end{frame}
%%%%%%%%%%%%%%%%%%%%%%%%%%%%%%%%%%%%%%%%%%%%%%%%%%%%%%%%%%%%%%%%%%%%%%%%%%%%%%%%
\begin{frame}
\frametitle{Objectives}
\label{sec-2}

  In this tutorial we will learn, \pause

\begin{itemize}
\item Create matrices using data.\pause
\item Create matrices from lists.\pause
\item Do basic matrix operations like addition, multiplication.
\end{itemize}
\end{frame}
%%%%%%%%%%%%%%%%%%%%%%%%%%%%%%%%%%%%%%%%%%%%%%%%%%%%%%%%%%%%%%%%%%%%%%%%%%%%%%%%
\begin{frame}
\frametitle{Objectives contd..}
\label{sec-2.2}
\begin{itemize}
\item Perform operations to find out the
	\begin{itemize}
	\item inverse of a matrix\pause
	\item determinant of a matrix\pause
	\item eigen values and eigen vectors of a matrix\pause
	\item norm of a matrix\pause
	\item singular value decomposition of a matrix.
	\end{itemize}
\end{itemize}
\end{frame}
%%%%%%%%%%%%%%%%%%%%%%%%%%%%%%%%%%%%%%%%%%%%%%%%%%%%%%%%%%%%%%%%%%%%%%%%%%%%%%%%
\begin{frame}
\frametitle{System Specifications}\pause
\begin{itemize}
\item Slackware Linux 14.1\pause
\item \texttt{Python 2.7.5} \pause
\item \texttt{IPython 2.3.0}
\end{itemize}
\end{frame}
%%%%%%%%%%%%%%%%%%%%%%%%%%%%%%%%%%%%%%%%%%%%%%%%%%%%%%%%%%%%%%%%%%%%%%%%%%%%%%%%
\begin{frame}
\frametitle{Pre-requisites}
\label{sec-3}

  Spoken tutorial on -\pause

\begin{itemize}
\item Getting started with Lists.\pause
\item Getting started with Arrays.\pause
\item Accessing parts of Arrays.
\end{itemize}
\end{frame}
%%%%%%%%%%%%%%%%%%%%%%%%%%%%%%%%%%%%%%%%%%%%%%%%%%%%%%%%%%%%%%%%%%%%%%%%%%%%%%%%
\begin{frame}
\frametitle{Exercise 1}
\label{sec-4}


\begin{itemize}
\item Create a two dimensional matrix \texttt{m3} of order \texttt{(2, 4)} with
   elements \\ \texttt{5, 6, 7, 8, 9, 10, 11, 12}.
\end{itemize}
\end{frame}
%%%%%%%%%%%%%%%%%%%%%%%%%%%%%%%%%%%%%%%%%%%%%%%%%%%%%%%%%%%%%%%%%%%%%%%%%%%%%%%%
\begin{frame}
\frametitle{Recall from \texttt{array}}
\label{sec-5.1}

  The following functions can also be used with matrices\pause

\begin{itemize}
\item \texttt{identity(n)}
	\begin{itemize}
	\item creates an identity matrix of order \texttt{nXn}\pause
	\end{itemize}
\item \texttt{zeros((m,n))}
	\begin{itemize}
	\item creates a matrix of order \texttt{mXn} with 0's
	\end{itemize}
\end{itemize}
\end{frame}
%%%%%%%%%%%%%%%%%%%%%%%%%%%%%%%%%%%%%%%%%%%%%%%%%%%%%%%%%%%%%%%%%%%%%%%%%%%%%%%%
\begin{frame}
\frametitle{Recall from \texttt{array}}
\label{sec-5.2}

\begin{itemize}
\item \texttt{zeros\_like(A)}
	\begin{itemize}
	\item creates a matrix with 0's similar to the shape of matrix \texttt{A}\pause
	\end{itemize}
\item \texttt{ones((m,n))}
	\begin{itemize}
	\item creates a matrix of order \texttt{mXn} with 1's\pause
	\end{itemize}
\item \texttt{ones\_like(A)}
	\begin{itemize}
	\item creates a matrix with 1's similar to the shape of matrix \texttt{A}
	\end{itemize}
\end{itemize}
\end{frame}
%%%%%%%%%%%%%%%%%%%%%%%%%%%%%%%%%%%%%%%%%%%%%%%%%%%%%%%%%%%%%%%%%%%%%%%%%%%%%%%%
\begin{frame}[fragile]
\frametitle{Exercise 2 : Frobenius norm \& inverse}
\label{sec-6}
\begin{itemize}
\item Frobenius norm is defined as,
	\begin{itemize}
	\item $||A||_F = \sqrt{\sum_{i,j}|a_{i,j}|^2}$\pause
	\end{itemize}
\begin{itemize}
\item Find out the Frobenius norm of inverse of a \texttt{4 X 4} matrix.
\end{itemize}
\texttt{m5=arange(1,17).reshape(4,4)}\pause
\item Inverse of A,
	\begin{itemize}
	\item $A^{-1} = inv(A)$
	\end{itemize}
\end{itemize}
\end{frame}
%%%%%%%%%%%%%%%%%%%%%%%%%%%%%%%%%%%%%%%%%%%%%%%%%%%%%%%%%%%%%%%%%%%%%%%%%%%%%%%%
\begin{frame}[fragile]
\frametitle{Exercise 3 : Infinity norm}
\label{sec-7}
\begin{itemize}
\item Infinity norm is defined as,
       $max(\sum_{i} |a_{i}|^2)$
\end{itemize}\pause
\begin{itemize}
\item Find the infinity norm of the matrix \texttt{im5}
\end{itemize}

\end{frame}
%%%%%%%%%%%%%%%%%%%%%%%%%%%%%%%%%%%%%%%%%%%%%%%%%%%%%%%%%%%%%%%%%%%%%%%%%%%%%%%%
\begin{frame}[fragile]
\frametitle{\texttt{norm()} method}
\label{sec-8}


\begin{itemize}
\item Frobenius norm\\
     \texttt{norm(im5)}\pause
\item Infinity norm\\
     \texttt{norm(im5, ord=inf)}
\end{itemize}
\end{frame}
%%%%%%%%%%%%%%%%%%%%%%%%%%%%%%%%%%%%%%%%%%%%%%%%%%%%%%%%%%%%%%%%%%%%%%%%%%%%%%%%
\begin{frame}
\frametitle{eigen values \& eigen vectors}
\label{sec-9}

  eigen values and eigen vectors

\begin{itemize}
\item \texttt{eig()}\pause
\end{itemize}
 
  Only eigen values

\begin{itemize}
\item \texttt{eigvals()}
\end{itemize}
\end{frame}
%%%%%%%%%%%%%%%%%%%%%%%%%%%%%%%%%%%%%%%%%%%%%%%%%%%%%%%%%%%%%%%%%%%%%%%%%%%%%%%%
\begin{frame}[fragile]
\frametitle{Singular Value Decomposition \texttt{svd}}
\label{sec-10.1}
   {\small  $M = U \Sigma V^*$}\pause
\begin{itemize}
\item U, an \texttt{mXm} unitary matrix over K.\pause
\item $\Sigma$, 
		 an \texttt{mXn} diagonal matrix with non-negative real numbers on diagonal.
\end{itemize}
\end{frame}
%%%%%%%%%%%%%%%%%%%%%%%%%%%%%%%%%%%%%%%%%%%%%%%%%%%%%%%%%%%%%%%%%%%%%%%%%%%%%%%%
\begin{frame}[fragile]
\frametitle{Singular Value Decomposition \texttt{svd} contd..}
\label{sec-10.2}
 
\begin{itemize}

\item $V^*$,
         an \texttt{nXn} unitary matrix over K, denotes the conjugate transpose of V.\pause
\item SVD of matrix \texttt{m5} can be found out as, \\ \texttt{svd(m5)}
\end{itemize}

\end{frame}
%%%%%%%%%%%%%%%%%%%%%%%%%%%%%%%%%%%%%%%%%%%%%%%%%%%%%%%%%%%%%%%%%%%%%%%%%%%%%%%%
\begin{frame}
\frametitle{Summary}
\label{sec-11.1}

  In this tutorial, we have learnt to - 
  
\begin{itemize}
\item Create matrices using arrays.
\item Add and multiply the elements of matrix. 
\item Find out the inverse of a matrix, using the function \texttt{inv()}. 
\end{itemize}
\end{frame}
%%%%%%%%%%%%%%%%%%%%%%%%%%%%%%%%%%%%%%%%%%%%%%%%%%%%%%%%%%%%%%%%%%%%%%%%%%%%%%%%
\begin{frame}
\frametitle{Summary contd..}
\label{sec-11.2}
\begin{itemize}
\item Use the function \texttt{det()} to find the determinant of a matrix. 
\item Calculate the norm of a matrix using the for loop and also using 
    the function \texttt{norm()}.
\end{itemize}
\end{frame}
%%%%%%%%%%%%%%%%%%%%%%%%%%%%%%%%%%%%%%%%%%%%%%%%%%%%%%%%%%%%%%%%%%%%%%%%%%%%%%%%
\begin{frame}
\frametitle{Summary contd..}
\label{sec-11.3}
\begin{itemize}
\item Find out the eigen vectors and eigen values of a matrix, using 
    functions \texttt{eig()} and \texttt{eigvals()}. 
\item Calculate singular value decomposition(SVD) of a matrix using the 
    function \texttt{svd()}.
\end{itemize}
\end{frame}
%%%%%%%%%%%%%%%%%%%%%%%%%%%%%%%%%%%%%%%%%%%%%%%%%%%%%%%%%%%%%%%%%%%%%%%%%%%%%%%%
\begin{frame}
\frametitle{Evaluation}
\label{sec-12.1}

\begin{enumerate}
\item \texttt{A} and \texttt{B} are two array objects. Element wise multiplication in
   matrices are done by,\pause
\begin{itemize}
\item \texttt{A * B}
\item \texttt{multiply(A, B)}
\item \texttt{dot(A, B)}
\item \texttt{element\_multiply(A,B)}\pause
\end{itemize}
\seti
\end{enumerate}
\end{frame}
%%%%%%%%%%%%%%%%%%%%%%%%%%%%%%%%%%%%%%%%%%%%%%%%%%%%%%%%%%%%%%%%%%%%%%%%%%%%%%%%
\begin{frame}
\frametitle{Evaluation contd..}
\label{sec-12.2}

\begin{enumerate}
\conti
\item \texttt{eig(A)[ 1 ]} and \texttt{eigvals(A)} are the same. (True of False)\pause
\item \texttt{norm(A, ord = 'fro')} is the same as \texttt{norm(A)} (True of False)
\end{enumerate}
\end{frame}
%%%%%%%%%%%%%%%%%%%%%%%%%%%%%%%%%%%%%%%%%%%%%%%%%%%%%%%%%%%%%%%%%%%%%%%%%%%%%%%%
\begin{frame}
\frametitle{Solutions}
\label{sec-13}


\begin{enumerate}
\item \texttt{A * B}\pause
\vspace{12pt}
\item False\pause
\vspace{12pt}
\item True
\end{enumerate}
\end{frame}
%%%%%%%%%%%%%%%%%%%%%%%%%%%%%%%%%%%%%%%%%%%%%%%%%%%%%%%%%%%%%%%%%%%%%%%%%%%%%%%%
\begin{frame}
\frametitle{FOSSEE}
{\color{blue}Free and Open-source Software for \\Science and Engineering Education} \\
\begin{itemize}
\item Goal - enabling all to use open source software tools
\item For more details, please visit {\color{blue}\url{http://fossee.in/}}
\end{itemize}
\end{frame}
%%%%%%%%%%%%%%%%%%%%%%%%%%%%%%%%%%%%%%%%%%%%%%%%%%%%%%%%%%%%%%%%%%%%%%%%%%%%%%%%
\begin{frame}
\frametitle{About the Spoken Tutorial Project}
\begin{itemize}
\item Watch the video available at {\color{blue}http://spoken-tutorial.org /What\_is\_a\_Spoken\_Tutorial}
\item It summarises the Spoken Tutorial project \pause
\item If you do not have good bandwidth, you can download and watch it
\end{itemize}
\end{frame}
%%%%%%%%%%%%%%%%%%%%%%%%%%%%%%%%%%%%%%%%%%%%%%%%%%%%%%%%%%%%%%%%%%%%%%%%%%%%%%%%
\begin{frame}
\frametitle{Spoken Tutorial Workshops}The Spoken Tutorial Project Team 
\begin{itemize}
\item Conducts workshops using spoken tutorials 
\item Gives certificates to those who pass an online test 
\item For more details, please write to \\ \hspace {0.5cm}{\color{blue}contact@spoken-tutorial.org}
\end{itemize}
\end{frame}
%%%%%%%%%%%%%%%%%%%%%%%%%%%%%%%%%%%%%%%%%%%%%%%%%%%%%%%%%%%%%%%%%%%%%%%%%%%%%%%%
\begin{frame}
\frametitle{Acknowledgements}
\begin{itemize}
\item Spoken Tutorial Project is a part of the Talk to a Teacher  project 
\item It is supported by the National Mission on Education through  ICT, MHRD, Government of India 
\item More information on this Mission is available at: \\{\color{blue}http://spoken-tutorial.org /NMEICT-Intro}
\end{itemize}
\end{frame}
%%%%%%%%%%%%%%%%%%%%%%%%%%%%%%%%%%%%%%%%%%%%%%%%%%%%%%%%%%%%%%%%%%%%%%%%%%%%%%%%
\begin{frame}

  \begin{block}{}
  \begin{center}
  \textcolor{blue}{\Large THANK YOU!} 
  \end{center}
  \end{block}
\begin{block}{}
  \begin{center}
    For more Information, visit our website\\
    \url{http://fossee.in/}
  \end{center}  
  \end{block}
\end{frame}
%%%%%%%%%%%%%%%%%%%%%%%%%%%%%%%%%%%%%%%%%%%%%%%%%%%%%%%%%%%%%%%%%%%%%%%%%%%%%%%%
\end{document}
